\section{Assignment 3}
\subsection{Introduction and objectives}
This assignment extends the damped wave simulator from parts 1 and 2 with
color. The underlying physics, discretisation, and MPI/OpenMP structure are
unchanged: three MPI ranks each run one independent scenario (a single
impulse, two impulses starting together, two impulses with the second
delayed to $N/7$), and each rank parallelises its own $M \times M$ stencil
update with OpenMP, exactly as in parts 1 and 2. What changes is the output
stage: instead of a grayscale PGM frame, each timestep's matrix is sent to
the GPU, mapped to an RGB triplet per pixel by a CUDA kernel, and saved as a
P3 (ASCII) PPM file. Bright, saturated colors mark the wave's extrema,
colors fade gradually along the rising/falling fronts, and near-zero, calm
regions are left white. The scalar matrix driving the simulation is never
modified by this process -- only the saved frame is colored.

Simulation parameters (student 352165), unchanged from parts 1--2:
\begin{center}
    \begin{tabular}{|c|c|c|c|c|}
    \hline
    \textbf{$\gamma$} & \textbf{$c$} & \textbf{$(i_0,j_0)$} & \textbf{$u(i_0,j_0)$} & \\
    \hline
    0.067 & 0.59 & $(M/2,M/2)$ & 56 & \\
    \hline
    \end{tabular}
\end{center}
sim2's second wave: $(2M/3,\,M/4)$, amplitude $-35$, $t_{\text{start}}=0$.
sim3's second wave: $(M/3,\,M/3)$, amplitude $58$, $t_{\text{start}}=N/7$.

\subsection{Implementation}
The project was rewritten from an earlier draft to remove unused
configurability (most physical parameters were assignment-mandated
constants, not things ever actually varied at runtime) and to make the
CUDA path measurably faster, not just correct.

\paragraph{CUDA colorization.} \texttt{colorize\_kernel} runs one GPU thread
per pixel. Each thread normalizes its value by a fixed scale (58, the
largest of the three configured amplitudes), applies a $\sqrt{|\cdot|}$
contrast boost so mid-amplitude fronts remain visually distinguishable
rather than washed out, and fades from white toward red/orange (positive)
or blue (negative) as magnitude grows past a small \texttt{zero\_band}
threshold -- satisfying the assignment's requirement for bright extrema,
graded fronts, and white calm regions.

\paragraph{Performance choices.} Three changes distinguish this
implementation from a direct "make it work" port: (1) the CUDA colorizer
context (device buffers, host buffers) is allocated once per scenario and
reused for all 300 frames, not re-allocated per frame; (2) the host-side RGB
buffer uses pinned (\texttt{cudaHostAlloc}) memory rather than plain
\texttt{malloc}, enabling faster DMA transfer; (3) the ASCII PPM writer
builds the entire frame's text in one pre-allocated buffer and issues a
single \texttt{fwrite}, instead of one \texttt{fprintf} call per pixel
value -- the latter was identified as the likely dominant CPU-side cost in
early testing, mirroring the \texttt{fwrite}-per-pixel bottleneck that part
1's VTune analysis found in the grayscale writer.

\paragraph{Build.} No Makefile is used. \texttt{assignment\_3.c} is
compiled and linked by \texttt{mpicc} (which already knows its own MPI
include/library paths); \texttt{wave\_color\_cuda.cu} is compiled separately
by \texttt{nvcc}; the two object files are linked together by \texttt{mpicc},
given only the CUDA runtime's library path.

\subsubsection{Engineering journey: errors and fixes}
Several real build- and tool-level errors were hit and resolved during
development; they are recorded here since diagnosing them shaped some of
the final design choices above.

\begin{itemize}
\item \textbf{\texttt{nvcc fatal: Unknown option '-Wl,-rpath'}.} Caused by
  piping \texttt{mpicc --showme:link}'s output directly into \texttt{nvcc}:
  that output includes \texttt{gcc}-style linker passthrough flags which
  \texttt{nvcc}'s argument parser does not accept the same way. Fixed by
  letting \texttt{mpicc} do all MPI-aware compiling and linking itself
  (see Build, above), so \texttt{nvcc} never has to see an MPI flag at all.
\item \textbf{Nsight Compute: ``option is ambiguous, matches
  `app-replay-buffer', `cache-control', \ldots''.} Caused by copying
  Nsight Systems' \texttt{-{}-} convention (``everything after this is the
  target executable'') into the Nsight Compute invocation. Nsight Compute
  does not use that separator; the stray \texttt{-{}-} was parsed as an
  option name of zero length, which is a valid (ambiguous) prefix of every
  long option Nsight Compute knows. Fixed by passing the target executable
  directly, with no separator.
\item \textbf{Nsight Compute: \texttt{ERR\_NVGPUCTRPERM}.} Once the above
  was fixed, Nsight Compute connected to all three ranks but was refused
  access to the GPU's performance counters -- a cluster-wide, admin-level
  restriction, not something fixable from a student account. Nsight Compute
  was dropped from the workflow; Nsight Systems (which does not need this
  permission) and the program's own internal timers cover the same ground
  for this report.
\item \textbf{Silent CSV column misalignment.} The sweep script's CSV
  headers for the thread- and block-size sweeps were written with one
  column name fewer than the number of fields actually being written per
  row. \texttt{pandas} does not error on this: it silently treats the first
  field as an unlabeled row index and shifts every header name one column
  over, so every plotted value was reading the wrong field. Caught by
  cross-checking the parsed DataFrame against the raw CSV text rather than
  trusting that ``the script ran without error'' meant the output was
  correct; fixed by correcting the header line to match the true column
  count.
\end{itemize}

\subsection{Testing and Results}
\label{sec:testing3}
Unless noted otherwise, all sweeps use \texttt{--steps 50} (fewer than the
production run's 300, to keep the whole sweep inside one job) and were run
on the \texttt{edu\_a40} partition (2 nodes, 48 listed cores, 4$\times$A40
GPUs per node; whether ``48'' is physical or logical cores -- \texttt{edu\_sapphire}
lists 48 but schedules 96 -- was still being confirmed as of this writing,
see the open item below).

\subsubsection{Sweep of grid size}
\begin{figure}[H]
    \centering
    \includegraphics[width=0.9\textwidth]{../src/Davide/Final/plots/size_sweep_execution_time.png}
    \caption{Execution time and its GPU-colorize/OpenMP-compute split, varying grid size $M$, 4 OpenMP threads, 3 concurrent ranks sharing one GPU.}
    \label{fig:size_sweep3}
\end{figure}
Execution time grows super-linearly with $M$ and approaches ideal
$O(M^2)$ scaling as $M$ increases: the $128\to256$ doubling costs only
$2.1\times$ (fixed per-run overhead still dominates at this scale), while
the $2048\to4096$ doubling costs $4.1\times$, essentially the ideal
quadratic ratio. The right panel makes the more important point plainly:
GPU colorize-and-write time dominates total execution time throughout, and
increasingly so as $M$ grows -- at $M=4096$ it is 79.3s of 81.1s total
($\sim$98\%), against OpenMP's \texttt{compute\_next} at 1.8s
($\sim$2\%). Both phases are $O(M^2)$, so this ratio is essentially
constant across the whole sweep (also $\sim$98\%/2\% at $M=2048$) --
coloring and writing an ASCII PPM frame is simply far more expensive per
pixel than the finite-difference update itself.

\subsubsection{Sweep of OpenMP thread count}
\begin{figure}[H]
    \centering
    \includegraphics[width=0.9\textwidth]{../src/Davide/Final/plots/thread_sweep_speedup_efficiency.png}
    \caption{Speedup and parallel efficiency vs. OpenMP thread count at $M=4096$, both for the full pipeline and for \texttt{compute\_next} in isolation.}
    \label{fig:thread_sweep3}
\end{figure}
Both curves are effectively flat, never exceeding roughly $1.2\times$
speedup even at 16 threads. This is not simply because $M$ was too small
(a first attempt at $M=1024$ showed the same flatness, motivating this
larger $M=4096$ rerun) -- both timed phases are $O(M^2)$, so their ratio,
and therefore how visible \texttt{compute\_next}'s own scaling is against
the dominant GPU cost, does not change with $M$. Looking at
\texttt{compute\_next}'s own time in isolation (the ``compute-only''
series) shows real but small speedup, saturating at $\sim$1.15--1.2$\times$
by 2 threads with no further gain through 16. This is consistent with the
5-point stencil update being memory-bandwidth-bound rather than
compute-bound: each output point costs 4 loads, 1 store, and a handful of
FLOPs, so a small number of threads is enough to saturate available memory
bandwidth, after which adding threads only adds contention for it. This
differs from part 1's OpenMP scaling results (which showed real speedup at
large thread counts) mainly because part 1's benchmark used $M$ up to
15000, at a scale where the total work is large enough to be interesting on
its own, independent of any GPU-bound stage.

\subsubsection{Sweep of CUDA block size}
\begin{figure}[H]
    \centering
    \includegraphics[width=0.75\textwidth]{../src/Davide/Final/plots/blocksize_sweep_kernel_time.png}
    \caption{GPU colorize-and-write time vs. CUDA threads-per-block, $M=4096$.}
    \label{fig:blocksize_sweep3}
\end{figure}
Both very small (32) and the maximum (1024) threads-per-block measurably
underperform a mid-range choice: 32 costs $\sim$88s on average, 1024
$\sim$88s, against a minimum of $\sim$78s around 128 threads/block --
roughly a 12\% swing. The same qualitative shape (worse at the extremes,
best in the low hundreds) appeared independently in an earlier sweep at
$M=1024$, which strengthens confidence this is a genuine occupancy effect
rather than run-to-run noise, though the curve itself is somewhat jagged
(plausible given 3 ranks concurrently contending for one physical GPU) --
averaging repeated runs per block size would sharpen it further if needed.

\subsubsection{GPU contention: solo vs. shared \textit{(pending rerun)}}
To separate genuine per-frame cost from contention for the one GPU all
three ranks share, \texttt{assignment\_3.c} was extended to accept an
explicit single-rank ``solo'' mode (previously rejected outright, since the
program assumed exactly the 3-scenario production layout), and
\texttt{sweep\_assignment3.sh} now runs sim1's workload both alone
(1 rank, exclusive GPU) and as part of the normal 3-rank run, at identical
parameters. \textit{Results pending a rerun of the sweep script; expected
figure: \texttt{contention\_sweep\_solo\_vs\_shared.png}.}

\subsubsection{Nsight Systems timeline \textit{(pending analysis)}}
\texttt{job\_assignment3.sh} produces one \texttt{.nsys-rep} timeline per
rank (\texttt{--trace=cuda,openmp,osrt}), viewed locally in the Nsight
Systems GUI. \textit{Not yet analyzed in this report; Nsight Compute
kernel-level metrics (occupancy, memory throughput) were attempted but are
unavailable on this cluster --} \texttt{ERR\_NVGPUCTRPERM} \textit{--
so the timeline view is the primary GPU-side profiling evidence for this
assignment.}

\subsubsection{Physical vs. logical core count on \texttt{edu\_a40}}
The HPC@PoliTO guide lists 48 cores for both \texttt{edu\_sapphire} and
\texttt{edu\_a40}, but \texttt{edu\_sapphire} was observed to actually
schedule 96. \texttt{sinfo -p edu\_a40 -N -o "\%N \%c \%X \%Y \%Z"} confirms
the same holds for \texttt{edu\_a40}: 2 sockets $\times$ 24 cores $\times$ 2
threads/core = 96 logical cores per node (SMT), not 48 -- the guide's figure
is physical cores only. The thread sweep in Figure~\ref{fig:thread_sweep3}
was capped at 16 threads accordingly; \texttt{sweep\_assignment3.sh} now
requests \texttt{--cpus-per-task=32} (the true per-rank ceiling, $96/3$
ranks) and tests up to 32 threads. \textit{Rerun pending; expectation, per
the memory-bandwidth-bound argument above, is that the plateau already
visible by 2 threads continues unchanged through 32.}
